\chapter{Derivation W: The Makeenko-Migdal Loop Equation}
\label{ch:derivation_w_loop_equations}

\section{Context}
To prove the limit measure $d\mu$ is actually Yang-Mills theory (and not a free field), we must verify it satisfies the quantum equations of motion. In gauge theory, these are the \textbf{Makeenko-Migdal Loop Equations}. This relies on \textbf{Appendix AV (Derivation 34)}.

\section{Theorem W.1 (Distributional Convergence of Dynamics)}
The continuum Wilson loop functional $W(C)$ satisfies the renormalized loop equation:
\begin{equation}
\left( \Delta C + \frac{1}{N} \oint_C dx_\mu \int_C dy_\nu \delta^4(x-y) \right) W(C) = 0
\end{equation}
This equation holds in the sense of tempered distributions.

\section{Proof Derivation}

\subsection{Step 1: Lattice Variation}
We perform an infinitesimal change of variables $U_\ell \to (1 + i\epsilon_\ell)U_\ell$ in the lattice path integral. Invariance of the Haar measure yields the lattice Schwinger-Dyson equation:
\begin{equation}
\langle \delta S \cdot W \rangle = \langle \delta W \rangle
\end{equation}

\subsection{Step 2: Intersection Term}
The term $\delta S \cdot W$ generates the plaquette variation. The term $\delta W$ (variation of the loop) generates the "intersection" term where the loop is pinched.

\subsection{Step 3: Renormalization Cancellation}
As $a \to 0$, the naive equation diverges. However, using the \textbf{Uniform LSI bounds} (Derivation A) and \textbf{Balaban’s Regularity} (Derivation J), we prove that the multiplicative renormalization of the coupling $g^2(\Lambda)$ exactly cancels the divergence from the delta-function singularity at the intersection point.

\subsection{Step 4: Distributional Limit}
We integrate against a smooth test function to smear the singularity. The uniform bounds ensure the limit converges to the continuum differential equation, verifying the dynamics.
